\DocumentMetadata{
  pdfversion=2.0,
  pdfstandard=ua-2,
  lang=en-US,
  tagging=on,
  tagging-setup={math/setup=mathml-SE,tabsorder=structure}
}
\documentclass[11pt,letterpaper]{article}
\usepackage[margin=0.68in,includefoot]{geometry}
\usepackage{newcomputermodern}
\usepackage{amsmath,amscd,mathtools}
\usepackage{tikz,adjustbox}
\providecommand{\square}{\mdlgwhtsquare}
\usepackage[hidelinks]{hyperref}
\hypersetup{
  pdftitle={Numerical Analysis PhD Exam, May 1999},
  pdfauthor={University of Florida Department of Mathematics},
  pdfsubject={Graduate mathematics examination},
  pdfdisplaydoctitle=true,
  bookmarksopen=true
}
\setcounter{secnumdepth}{0}
\setlength{\parindent}{0pt}
\setlength{\parskip}{0.28em}
\setlength{\emergencystretch}{3em}
\setlength{\leftmargini}{2.15em}
\setlength{\leftmarginii}{2.65em}
\setlength{\leftmarginiii}{3em}
\pagestyle{empty}
\raggedbottom
\allowdisplaybreaks
\NewTaggingSocketPlug{block/list/label}{examproblem}{%
  \tagstructbegin{tag=Lbl}\tagstructend%
  \tagstructbegin{tag=\LBody}%
  \tagstructbegin{tag=\ProblemHeadingTag,title-o={\ProblemHeadingText}}%
  \tagmcbegin{tag=\ProblemHeadingTag,actualtext={\ProblemHeadingText}}%
  #2%
  \tagmcend%
  \tagstructend%
}
\newenvironment{examproblems}[2]{%
  \def\ProblemHeadingTag{#2}%
  \AssignTaggingSocketPlug{block/list/label}{examproblem}%
  \begin{enumerate}%
  \setcounter{enumi}{#1}%
  \renewcommand{\labelenumi}{\textbf{\theenumi.}}%
  \setlength{\topsep}{0.35em}%
  \setlength{\partopsep}{0pt}%
  \setlength{\itemsep}{0.48em}%
  \setlength{\parsep}{0.18em}%
}{\end{enumerate}}
\newenvironment{examsubparts}{%
  \AssignTaggingSocketPlug{block/list/label}{default}%
  \begin{enumerate}%
  \setlength{\topsep}{0.18em}%
  \setlength{\partopsep}{0pt}%
  \setlength{\itemsep}{0.16em}%
  \setlength{\parsep}{0.1em}%
}{\end{enumerate}}
\newenvironment{examinstructions}{%
  \par\begingroup\itshape
}{%
  \par\endgroup\vspace{0.18em}
}
\makeatletter
\renewcommand\subsection{\@startsection{subsection}{2}{0pt}%
  {-1.25ex plus -.25ex minus -.1ex}%
  {0.55ex plus .1ex}%
  {\normalfont\large\bfseries}}
\renewcommand{\@maketitle}{%
  \newpage\null\vskip -1em%
  {\footnotesize\bfseries
    \href{https://www.ufl.edu/}{University of Florida}, \href{https://math.ufl.edu/}{Department of Mathematics}\par}%
  \vskip -0.5em%
  \tagstructbegin{tag=H1,title={Numerical Analysis PhD Exam, May 1999}}%
  \tagpdfparaOff%
  \tagmcbegin{tag=H1}%
    {\LARGE\bfseries\href{https://gma.math.ufl.edu/exams/numerical-analysis-phd/}{Numerical Analysis PhD Exam}, May 1999\par}%
  \tagmcend%
  \tagpdfparaOn%
  \tagstructend%
  \vskip 0.25em%
  \tagmcbegin{artifact}%
    \hrule height 1pt%
  \tagmcend%
  \vskip 1em%
}
\makeatother
\title{Numerical Analysis PhD Exam, May 1999}
\author{}
\date{}
\begin{document}
\maketitle
\thispagestyle{empty}
\begin{examinstructions}
Do any 8 of the following 10 problems.
\end{examinstructions}
\begin{examproblems}{0}{H2}
\def\ProblemHeadingText{Problem 1.}
\item
This problem has six parts.
\begin{examsubparts}
\item[\textnormal{(1)}]
Give a careful statement and proof of the singular value decomposition.
\item[\textnormal{(2)}]
Give a careful statement of the \(QR\) factorization theorem.
\item[\textnormal{(3)}]
Indicate how the Householder transformations are used in the proof of the \(QR\) factorization theorem.
\item[\textnormal{(4)}]
Describe the \(QR\) Iteration Algorithm with shift~\(\mu\).
\item[\textnormal{(5)}]
Describe how the \(QR\) Iteration is used to compute the \(SVD\).
\item[\textnormal{(6)}]
How can the singular values be used to compute the condition number of a matrix.
\end{examsubparts}
\def\ProblemHeadingText{Problem 2.}
\item
Give a careful proof of the following Theorem: If \(p\geq 1\),

\par
1. \(\mathbf A\in\mathbb R^{n\times n}\) is invertible, 2. \(\mathbf b\) and \(\Delta\mathbf b\) are vectors in \(\mathbb R^n\), and 3. \(\mathbf x\) and \(\Delta\mathbf x\) are vectors in \(\mathbb R^n\) and are solutions to \(\mathbf A\mathbf x=\mathbf b\), and \(\mathbf A\Delta\mathbf x=\Delta\mathbf b\), respectively,

\par
then
\[
\frac{\|\Delta\mathbf x\|_p}{\|\mathbf x\|_p}\leq \kappa_p(\mathbf A)\frac{\|\Delta\mathbf b\|_p}{\|\mathbf b\|_p}.
\]
\def\ProblemHeadingText{Problem 3.}
\item
This problem has three parts.
\begin{examsubparts}
\item[\textnormal{(1)}]
Describe how Jacobi rotations can be used to diagonalize a real symmetric matrix.
\item[\textnormal{(2)}]
Give the error estimate for a single Jacobi rotation.
\item[\textnormal{(3)}]
Give the error estimates for repeated Jacobi rotations.
\end{examsubparts}
\def\ProblemHeadingText{Problem 4.}
\item
This problem has two parts.
\begin{examsubparts}
\item[\textnormal{(1)}]
Write out the system of linear equations to be solved to find the periodic spline interpolation for the given data points \((x_0,y_0),(x_1,y_1),\ldots,(x_n,y_n)\), where \(x_0<x_1<\cdots<x_n\) and \(y_n=y_0\).
\item[\textnormal{(2)}]
What method would you recommend for solving the system of linear equations. Explain!!
\end{examsubparts}
\def\ProblemHeadingText{Problem 5.}
\item
This problem has three parts.
\begin{examsubparts}
\item[\textnormal{(1)}]
Give a careful definition of the Householder transformation.
\item[\textnormal{(2)}]
Explain why the Householder transformation is preferred over the Gram–Schmidt orthogonalization process for creating an orthonormal basis. (Give an example if necessary.)
\item[\textnormal{(3)}]
Show how Householder transformations can be used to bidiagonalize any real matrix.
\end{examsubparts}
\def\ProblemHeadingText{Problem 6.}
\item
This problem has four parts.
\begin{examsubparts}
\item[\textnormal{(1)}]
Give a careful statement of the Pythagorean Theorem property for clamped cubic splines.
\item[\textnormal{(2)}]
Give a careful statement of the integral minimization property for clamped cubic splines.
\item[\textnormal{(3)}]
If \(s_n(x)\) denotes the clamped cubic spline interpolant for a function \(f(x)\) at the knots \((x_0,y_0),(x_1,y_1),\ldots,(x_n,y_n)\), then give a precise statement for the error between \(s_n(x)\) and \(f(x)\).
\item[\textnormal{(4)}]
Prove the error formula given in step (3) above.
\end{examsubparts}
\def\ProblemHeadingText{Problem 7.}
\item
This problem has three parts.
\begin{examsubparts}
\item[\textnormal{(1)}]
If \(p_n(x)=\frac{1}{2^n n!}\frac{d^n}{dx^n}(x^2-1)^n\), then show that \(\int_{-1}^1 x^m p_n(x)\,dx=0\) for all \(m<n\). (Hint: Think small values of~\(m\) and~\(n\), induction, and parts.)
\item[\textnormal{(2)}]
What can you conclude about the relationship between the polynomials \(p_n(x)\) and the Legendre polynomials?
\item[\textnormal{(3)}]
Explain the method of Gauss quadrature for the numerical approximation of the integral of a function.
\end{examsubparts}
\def\ProblemHeadingText{Problem 8.}
\item
This problem has three parts.
\begin{examsubparts}
\item[\textnormal{(1)}]
Discuss the ideas behind the linear shooting method for approximating the solution of a linear two point boundary value problem. In particular, explain how the Runge–Kutta method can be used in the approximation.
\item[\textnormal{(2)}]
Describe how the Galerkin method can be used to solve a 2 point boundary value problem on the interval \([0,1]\).
\item[\textnormal{(3)}]
Explain why splines have an advantage over \(\sin(kx)\) as choices of test functions for the Galerkin method.
\end{examsubparts}
\def\ProblemHeadingText{Problem 9.}
\item
This problem has three parts.
\begin{examsubparts}
\item[\textnormal{(1)}]
Give a careful statement of the Contraction Mapping Theorem.
\item[\textnormal{(2)}]
Outline a proof of the Contraction Mapping Theorem.
\item[\textnormal{(3)}]
Indicate how Aitken’s method can be used to accelerate convergence.
\end{examsubparts}
\def\ProblemHeadingText{Problem 10.}
\item
This problem has two parts.
\begin{examsubparts}
\item[\textnormal{(1)}]
Give a careful statement of the error formula for the trapezoidal rule for approximating an integral.
\item[\textnormal{(2)}]
Give a careful proof of the error formula for the trapezoidal rule.
\end{examsubparts}
\end{examproblems}
\end{document}
